[[1.1]] Compton scattering: a photon of wavelength $\lambda$ scatters on a particle of mass $m$ at rest. Find the new photon wavelength $\lambda^{\prime}$, provided the scattering angle is $\theta$.
[[2.1]] Verify that $\gamma^{\mu \dagger}=\gamma^0 \gamma^\mu \gamma^0$. Show that the adjoint spinor transforms as $\bar{\psi}^{\prime}\left(x^{\prime}\right)= \bar{\psi}(x) S^{-1}(\Lambda)$. Some covariant bilinears $\bar{\psi} \Gamma \psi$ correspond to observable quantities and therefore must be real. Show that this implies, for the $\Gamma$ matrices, the relation $\Gamma=\bar{\Gamma}$, where $\bar{\Gamma}=\gamma^0 \Gamma^{\dagger} \gamma^0$.
[[3.1]] Use the Dirac equation and its adjoint to derive a continuity equation. Show that the four-current is given by

$$
j^\mu=\bar{\psi} \gamma^\mu \psi .
$$


Is $j^0=\rho$ positive definite or not? Prove it.
[[4.1]] Prove that the following bilinear covariants transform, under a general Lorentz transformation $S(\Lambda)$, in the way specified below.

\begin{tabular}{rll}
$\bar{\psi} \psi$ & $\rightarrow \bar{\psi} \psi$ & \\
$\bar{\psi} \gamma^5 \psi$ & $\rightarrow \operatorname{det}(\Lambda) \bar{\psi} \gamma^5 \psi$ & \\
$\bar{\psi} \gamma^\mu \psi$ & $\rightarrow \Lambda_\nu^\mu \bar{\psi} \gamma^\nu \psi$ & scalar \\
$\bar{\psi} \gamma^\mu \gamma^5 \psi$ & $\rightarrow \operatorname{det}(\Lambda) \Lambda_\nu^\mu \bar{\psi} \gamma^\nu \gamma^5 \psi$ & vector \\
$\bar{\psi} \sigma^{\mu \nu} \psi$ & $\rightarrow \Lambda_\alpha^\mu \Lambda_\beta^\nu \bar{\psi} \sigma^{\alpha \beta} \psi$ & axial vector \\
& & antisymmetric tensor
\end{tabular}
[[5.1]] The projectors on the left- and right-components of Dirac spinors are given by $P_L=\frac{1-\gamma^5}{2}$ and $P_R=\frac{1+\gamma^5}{2}$ respectively. Use the properties of $\gamma^5$ to prove the following projector identities:

$$
P_L P_R=P_R P_L=0, \quad P_R P_R=P_R, \quad P_L P_L=P_L .
$$


Verify then that, given a spinorial solution to the massless Dirac equation, the spinors $P_{L, R} u_s(p)$ are eigenstates of the helicity operator

$$
\mathfrak{h}=\frac{1}{2|\mathbf{p}|}\left(\begin{array}{cc}
\sigma \cdot \mathbf{p} & 0 \\
0 & \sigma \cdot \mathbf{p}
\end{array}\right)
$$

with eigenvalues $\pm \frac{1}{2}$. Does this hold also for $m \neq 0$ ?